\def\tituloconcepto{Crecimiento exponencial: $2^n$ amplitudes}
% Si el capitulo se compila solo (sin main.tex), esta cabecera abre el documento.
% Antes de % Si el capitulo se compila solo (sin main.tex), esta cabecera abre el documento.
% Antes de % Si el capitulo se compila solo (sin main.tex), esta cabecera abre el documento.
% Antes de \input{cabecera} cada capitulo define \tituloconcepto.
\ifdefined\enmain\else
  \documentclass[11pt,a4paper]{article}
  \input{preambulo}
  \begin{document}
  \title{\tituloconcepto}
  \author{\autor}
  \date{\today}
  \maketitle
\fi
 cada capitulo define \tituloconcepto.
\ifdefined\enmain\else
  \documentclass[11pt,a4paper]{article}
  % Preambulo compartido por todos los documentos. Compilar con XeLaTeX/LuaLaTeX
% (o Tectonic). Si se usa Overleaf, elegir el compilador XeLaTeX.
\usepackage[spanish,es-noshorthands]{babel}
\usepackage{amsmath,amssymb}
\usepackage{booktabs,array}
\usepackage[a4paper,margin=2.5cm]{geometry}
\usepackage{xcolor}
\usepackage{listings}
\usepackage[most]{tcolorbox}
\usepackage{quantikz}
\usepackage[hidelinks]{hyperref}

% ---- Autor (editar a gusto) -------------------------------------------------
\newcommand{\autor}{Genesis Silesky}

% ---- Notacion de Dirac y atajos ---------------------------------------------
% quantikz ya define \ket y \bra; DeclareRobustCommand los redefine sin error.
\DeclareRobustCommand{\ket}[1]{\left|#1\right\rangle}
\DeclareRobustCommand{\bra}[1]{\left\langle#1\right|}
\DeclareRobustCommand{\braket}[2]{\left\langle#1\middle|#2\right\rangle}
\newcommand{\C}{\mathbb{C}}
\newcommand{\R}{\mathbb{R}}
\DeclareMathOperator{\Tr}{Tr}

% ---- Cajas de apoyo --------------------------------------------------------
% "concepto": definiciones previas que hacen falta para entender el capitulo.
\newtcolorbox{concepto}[1]{breakable,colback=blue!4,colframe=blue!55!black,
  fonttitle=\bfseries,title={Concepto previo: #1}}
% "nota": aclaraciones y advertencias.
\newtcolorbox{nota}[1]{breakable,colback=orange!6,colframe=orange!70!black,
  fonttitle=\bfseries,title={#1}}

% ---- Estilos de listados ----------------------------------------------------
\lstdefinestyle{py}{language=Python,basicstyle=\ttfamily\footnotesize,
  breaklines=true,postbreak=\mbox{\textcolor{gray}{$\hookrightarrow$}\space},
  numbers=left,numberstyle=\tiny\color{gray},frame=single,
  commentstyle=\color{green!40!black}\itshape,
  keywordstyle=\color{blue!70!black}\bfseries,
  stringstyle=\color{red!60!black},showstringspaces=false,tabsize=4,
  columns=fullflexible,keepspaces=true}
\lstdefinestyle{salida}{basicstyle=\ttfamily\footnotesize,breaklines=true,
  frame=leftline,backgroundcolor=\color{gray!8},columns=fullflexible,
  keepspaces=true}

  \begin{document}
  \title{\tituloconcepto}
  \author{\autor}
  \date{\today}
  \maketitle
\fi
 cada capitulo define \tituloconcepto.
\ifdefined\enmain\else
  \documentclass[11pt,a4paper]{article}
  % Preambulo compartido por todos los documentos. Compilar con XeLaTeX/LuaLaTeX
% (o Tectonic). Si se usa Overleaf, elegir el compilador XeLaTeX.
\usepackage[spanish,es-noshorthands]{babel}
\usepackage{amsmath,amssymb}
\usepackage{booktabs,array}
\usepackage[a4paper,margin=2.5cm]{geometry}
\usepackage{xcolor}
\usepackage{listings}
\usepackage[most]{tcolorbox}
\usepackage{quantikz}
\usepackage[hidelinks]{hyperref}

% ---- Autor (editar a gusto) -------------------------------------------------
\newcommand{\autor}{Genesis Silesky}

% ---- Notacion de Dirac y atajos ---------------------------------------------
% quantikz ya define \ket y \bra; DeclareRobustCommand los redefine sin error.
\DeclareRobustCommand{\ket}[1]{\left|#1\right\rangle}
\DeclareRobustCommand{\bra}[1]{\left\langle#1\right|}
\DeclareRobustCommand{\braket}[2]{\left\langle#1\middle|#2\right\rangle}
\newcommand{\C}{\mathbb{C}}
\newcommand{\R}{\mathbb{R}}
\DeclareMathOperator{\Tr}{Tr}

% ---- Cajas de apoyo --------------------------------------------------------
% "concepto": definiciones previas que hacen falta para entender el capitulo.
\newtcolorbox{concepto}[1]{breakable,colback=blue!4,colframe=blue!55!black,
  fonttitle=\bfseries,title={Concepto previo: #1}}
% "nota": aclaraciones y advertencias.
\newtcolorbox{nota}[1]{breakable,colback=orange!6,colframe=orange!70!black,
  fonttitle=\bfseries,title={#1}}

% ---- Estilos de listados ----------------------------------------------------
\lstdefinestyle{py}{language=Python,basicstyle=\ttfamily\footnotesize,
  breaklines=true,postbreak=\mbox{\textcolor{gray}{$\hookrightarrow$}\space},
  numbers=left,numberstyle=\tiny\color{gray},frame=single,
  commentstyle=\color{green!40!black}\itshape,
  keywordstyle=\color{blue!70!black}\bfseries,
  stringstyle=\color{red!60!black},showstringspaces=false,tabsize=4,
  columns=fullflexible,keepspaces=true}
\lstdefinestyle{salida}{basicstyle=\ttfamily\footnotesize,breaklines=true,
  frame=leftline,backgroundcolor=\color{gray!8},columns=fullflexible,
  keepspaces=true}

  \begin{document}
  \title{\tituloconcepto}
  \author{\autor}
  \date{\today}
  \maketitle
\fi


\section{Crecimiento exponencial: $2^n$ amplitudes}

\subsection{Objetivo}
Cuantificar por qué simular $n$ qubits en una computadora clásica requiere
memoria que crece como $2^n$.

\subsection{Conceptos previos}
\begin{concepto}{producto tensorial y potencias de 2}
Ver sección~\ref{sec:prerreq}. Idea clave: al combinar sistemas con
$\otimes$, las dimensiones se \emph{multiplican}: $2\cdot2\cdots2=2^n$. Además,
$1\,\mathrm{GiB}=2^{30}$ bytes.
\end{concepto}

\subsection{Fundamento teórico}
El espacio de $n$ qubits es $(\C^2)^{\otimes n}\cong\C^{2^n}$; un estado
general es
\begin{equation}
  \ket\psi=\sum_{x\in\{0,1\}^n}c_x\ket x,\qquad \sum_x|c_x|^2=1,
\end{equation}
con $2^n$ amplitudes complejas independientes (salvo normalización y fase
global). Con \texttt{complex128} cada amplitud ocupa $16$ bytes, de modo que
\begin{equation}
  \text{memoria}=16\cdot2^n\ \text{bytes}.
\end{equation}
Para $n=30$: $16\cdot2^{30}$ B $=16$ GiB; para $n=50$: $16\cdot2^{50}$ B
$=2^{24}$ GiB $=16$ PiB, y hay $\approx1.13\times10^{15}$ amplitudes.

\paragraph{Un caso concreto.} Aplicar $H$ a cada uno de $n$ qubits partiendo de
$\ket{0\cdots0}$ da
$H^{\otimes n}\ket{0}^{\otimes n}=\frac1{\sqrt{2^n}}\sum_x\ket x$. Con $n=4$:
$16$ amplitudes iguales a $1/4$.

\begin{nota}{Matiz importante}
Esto aplica a la simulación de circuitos \emph{generales} con el vector de
estado. Hay circuitos con estructura (p.\,ej.\ solo puertas Clifford, o con
poco entrelazamiento) que se simulan eficientemente con otros métodos. Tampoco
significa que un computador cuántico ``guarde'' $2^n$ números accesibles: al
medir solo se obtienen $n$ bits.
\end{nota}

\subsection{Código (completamente comentado)}
\lstinputlisting[style=py]{../../fuentes/03_crecimiento_exponencial.py}

\subsection{Explicación celda por celda}
\begin{description}
  \item[Celda 2] Recorre varios $n$, calcula $2^n$ y convierte a GiB
    dividiendo entre $2^{30}$. Los formatos \texttt{:>18,} alinean y ponen
    separadores de miles.
  \item[Celda 3] Construye el estado $H^{\otimes4}\ket{0000}$ y comprueba que
    hay $16$ amplitudes, todas iguales a $0.25$.
\end{description}

\subsection{Salida obtenida}
\lstinputlisting[style=salida]{../../fuentes/salidas/03_crecimiento_exponencial.txt}

\subsection{Interpretación}
Cada qubit extra \emph{duplica} la memoria. A $30$ qubits ya se necesitan
$16$~GiB; unos pocos más y ninguna laptop alcanza. Este es el motivo de que
simular sea difícil y, a la vez, parte de la motivación de usar hardware
cuántico.

\subsection{Ejercicios}
\begin{enumerate}
  \item ¿Cuántos qubits se pueden simular con 64~GiB de RAM?
  \item Repita el cálculo con \texttt{complex64} (8 bytes por amplitud).
  \item Verifique con \texttt{Statevector} que para $n=3$ las amplitudes valen
    $1/\sqrt8\approx0.354$.
\end{enumerate}

\ifdefined\enmain\else
  \end{document}
\fi

